% Part II > The proof. Reconceptualizes Part I's proof-fidelity criterion: a
% proof is an auditable graph of formal statements (the scaffold), understood
% when a reader can walk it without opening any proof. Order (Kyle,
% 2026-11-11): the back and forth on the statement (previous section), then on
% the proof, then how a proof decomposes into interpretable lemmas. Where formal proofs come from;
% decomposing a proof into interpretable lemmas. No running example: this is
% not a Lean tutorial (Kyle, 2026-11-11).
% talk-v2, 2026-11-11.
\section{Proof Fidelity}

% [2 min] Where formal proofs come from (built 2026-11-11): from the statement
% alone versus from a written proof. Sources: AlphaProof Nexus, arXiv
% 2605.22763 (9/353 open Erdos, 44/492 OEIS, a few hundred dollars each;
% research/02 E); FLT 11 days on Buzzard's project, ~6B output tokens
% (Anthropic, self-reported; research/01); CFSG 24 days from 715 papers and
% books (Axiom, self-reported); Meta textbook scaffold, arXiv 2604.03071;
% sphere packing on a human blueprint (Math, Inc.); token estimates for
% Navier-Stokes (~130B) and FLT (~6B) from AcadMathSci B.1 (different
% problems: a discovery versus a formalization). Draft, Sketch, Prove (Jiang et
% al., arXiv 2210.12283): sketches from informal proofs raise success on
% competition problems from 20.9% to 39.3% (abstract; Isabelle, not Lean).
% DeepSeek-Prover-V2 (arXiv 2504.21801): training starts by decomposing
% problems into subgoals.
\begin{frame}{Where formal proofs come from}
  \centering
  \begin{tikzpicture}[x=1cm, y=1cm,
      mini/.style={draw=black!30, fill=white, rounded corners=2pt, line width=0.5pt,
                   minimum height=0.6cm, inner sep=0.08cm, align=center, font=\sffamily\scriptsize},
      colhd/.style={anchor=north west, inner sep=0pt, font=\sffamily\small\bfseries},
      body/.style={anchor=north west, text width=6.5cm, align=flush left, inner sep=0pt, font=\sffamily\scriptsize}]
    \useasboundingbox (0,0.1) rectangle (14.3,-5.9);
    \draw[black!12, line width=0.5pt] (7.15,0.05) -- (7.15,-3.65);
    % ================= from the statement alone =================
    \node[colhd, text=tlPurple] at (0,0.05) {From the statement alone};
    \node[mini, draw=PennBlue!55, fill=PennBlue!5, text width=1.7cm] (a1) at (1.0,-0.85) {Statement};
    \node[mini, draw=PennGreen!55, fill=PennGreen!6, text width=1.7cm] (a2) at (5.2,-0.85) {Lean proof};
    \draw[flow=tlPurple] (a1) -- node[note, above] {machine search} (a2);
    \node[body] at (0,-1.45) {%
      \begin{itemize}\setlength{\itemsep}{2pt}
        \item AlphaProof Nexus: 9 of 353 open Erd\H{o}s problems and 44 of 492 OEIS conjectures, at a few hundred dollars each.
        \item Works best on short, self-contained problems, where no written proof exists to follow.
        \item Navier--Stokes, found from the statement: about 130B tokens.
      \end{itemize}};
    % ================= from a written proof =================
    \node[colhd, text=PennGreen] at (7.5,0.05) {From a written proof};
    \node[mini, draw=PennBlue!55, fill=PennBlue!5, text width=2.3cm] (b1) at (8.75,-0.85) {Statement + written proof};
    \node[mini, draw=PennGreen!55, fill=PennGreen!6, text width=1.7cm] (b2) at (13.2,-0.85) {Lean proof};
    \draw[flow=tlPurple] (b1) -- node[note, above] {fill in} (b2);
    \node[body] at (7.5,-1.45) {%
      \begin{itemize}\setlength{\itemsep}{2pt}
        \item Fermat's Last Theorem in 11 days, on Buzzard's project and the literature: about 6B tokens.
        \item The finite simple groups in 24 days, from 715 papers and books.
        \item A 500-page textbook: all 340 targets in about a week. Sphere packing: built on a human blueprint.
      \end{itemize}};
    % ================= controlled comparisons =================
    \draw[black!15, line width=0.4pt] (0,-3.8) -- (14.3,-3.8);
    \node[anchor=north west, text width=14.2cm, align=flush left, inner sep=0pt, font=\sffamily\scriptsize] at (0,-4.0)
      {\textbf{Like for like.} Sketches from informal proofs raised a prover's success on competition problems from 20.9\% to 39.3\%
       (Draft, Sketch, Prove). {\color{black!55}Token counts are estimates for different problems.}\par\smallskip
       \textbf{And the written proof supplies the lemmas:} its structure is what we decompose next.};
    \node[keybox, anchor=south] at (7.15,-5.85) {A written proof is still the cheapest route to a formal one.};
  \end{tikzpicture}
  \slidesources{Google DeepMind, AlphaProof Nexus, arXiv 2605.22763 \textbullet{} Anthropic, FLT (Sep 2026) \textbullet{} Axiom, CFSG (Oct 2026) \textbullet{} Meta, arXiv 2604.03071 \textbullet{} AcadMathSci, Oct 2026 (B.1) \textbullet{} Jiang et al., arXiv 2210.12283}
\end{frame}

% [3 min] Decomposing a proof into interpretable lemmas (built 2026-11-11).
% Diagram: a statement node T with dependencies A, B, C and one large proof
% expands into lemmas L1, L2 that inherit some of those dependencies (L1: A, B;
% L2: B), while T keeps C directly. Proofs are drawn as small dark blocks:
% checked by the kernel, never read. Evidence: Axiom rebuilt its CFSG plan
% when prerequisites were missing (research/01 s2); sphere packing's proof went
% from about 450k to about 200k lines (Math, Inc., self-reported); LeanArchitect
% extracts the graph from code but does not choose interpretable lemmas
% (research/02 C). Key box wording agreed with Kyle (2026-11-11).
\begin{frame}{Decomposing the proof into interpretable lemmas}
  \centering
  \begin{tikzpicture}[x=1cm, y=1cm,
      dnode/.style={draw=black!35, fill=black!5, rounded corners=3pt, line width=0.5pt,
                    minimum width=1.35cm, minimum height=0.5cm, inner sep=0.05cm, font=\sffamily\scriptsize, text=black!70},
      lnode/.style={draw=PennGold!75, fill=PennGold!12, rounded corners=3pt, line width=0.6pt,
                    minimum width=1.7cm, minimum height=0.55cm, inner sep=0.05cm, font=\sffamily\scriptsize},
      tnode/.style={draw=PennBlue!60, fill=PennBlue!8, rounded corners=3pt, line width=0.6pt,
                    minimum width=2.0cm, minimum height=0.6cm, inner sep=0.05cm, font=\sffamily\scriptsize\bfseries},
      pf/.style={fill=black!70, draw=none},
      colhd/.style={anchor=north west, inner sep=0pt, font=\sffamily\small\bfseries, text=black!75},
      body/.style={anchor=north west, text width=6.6cm, align=flush left, inner sep=0pt, font=\sffamily\scriptsize}]
    \useasboundingbox (0,0.1) rectangle (14.3,-5.9);
    % ================= before =================
    \node[anchor=north west, inner sep=0pt, font=\sffamily\scriptsize\bfseries, text=black!55] at (0,0.05) {Before: one statement, one opaque proof};
    \foreach \n/\x/\t in {a/0.85/A, b/2.6/B, c/4.35/C} \node[dnode] (b\n) at (\x,-0.6) {\t};
    \node[tnode] (bt) at (2.6,-1.6) {Theorem};
    \foreach \n in {a,b,c} \draw[dep] (b\n) -- (bt);
    \fill[pf] (1.85,-1.98) rectangle (3.35,-2.62);
    \node[font=\sffamily\tiny, text=white] at (2.6,-2.3) {one long proof};
    % ================= the move =================
    \draw[-Stealth, line width=1.2pt, draw=PennGold!85!black] (5.45,-1.45) -- (6.85,-1.45);
    \node[anchor=south, font=\sffamily\scriptsize\bfseries, text=PennGold!85!black] at (6.15,-1.4) {decompose};
    % ================= after =================
    \node[anchor=north west, inner sep=0pt, font=\sffamily\scriptsize\bfseries, text=black!55] at (7.3,0.05) {After: lemmas that inherit the dependencies};
    \foreach \n/\x/\t in {a/8.1/A, b/10.2/B, c/13.2/C} \node[dnode] (a\n) at (\x,-0.6) {\t};
    \node[lnode] (l1) at (8.6,-1.45) {Lemma 1};
    \node[lnode] (l2) at (11.0,-1.45) {Lemma 2};
    \node[tnode] (at) at (10.6,-2.3) {Theorem};
    \draw[dep] (aa) -- (l1); \draw[dep] (ab) -- (l1); \draw[dep] (ab) -- (l2);
    \draw[dep] (ac) -- (at);
    \draw[dep] (l1) -- (at); \draw[dep] (l2) -- (at);
    \foreach \n in {l1, l2, at} \fill[pf] ([xshift=-0.05cm, yshift=0.05cm]\n.south east) rectangle ++(-0.26,-0.18);
    % legend
    \node[anchor=north west, inner sep=0pt, font=\sffamily\tiny, text=black!60] at (7.3,-2.8)
      {\tikz\fill[black!5, draw=black!35] (0,0) rectangle (0.25,0.15);\;depended on (library or earlier nodes)\quad
       \tikz\fill[PennGold!12, draw=PennGold!75] (0,0) rectangle (0.25,0.15);\;new lemma\quad
       \tikz\fill[black!70] (0,0) rectangle (0.25,0.15);\;proof: kernel-checked, unread};
    % ================= when to split / who chooses =================
    \draw[black!15, line width=0.4pt] (0,-3.15) -- (14.3,-3.15);
    \node[colhd] at (0,-3.3) {When to split a node};
    \node[body, text=PennRed!85!black] at (0,-3.65)
      {\bfseries Proof fidelity: the lemma graph matches, not every line.};
    \node[body] at (0,-3.82) {%
      \begin{itemize}\setlength{\itemsep}{0pt}
        \item Its statement is not a building block the audience can hold.
        \item Its proof is long next to its statement: something is hidden.
        \item Only statements are read; the route inside does not matter.
      \end{itemize}};
    \node[colhd] at (7.3,-3.3) {Who chooses the lemmas};
    \node[body, text width=6.9cm] at (7.3,-3.7) {%
      \begin{itemize}\setlength{\itemsep}{1pt}
        \item Back and forth: the human proposes nodes; the machine fills them, fails, or proposes lemmas; the human keeps those that read clearly.
        \item Tools extract the graph (LeanArchitect); none yet chooses interpretable lemmas.
      \end{itemize}};
    \node[keybox, anchor=south] at (7.15,-6) {Lemmas comprise the signature of human perspective and proof authorship.};
  \end{tikzpicture}
  \slidesources{Zhu et al., LeanArchitect, ITP 2026 \textbullet{} research note 02 (C, D)}
\end{frame}

% [2 min] Closer for Part II (built 2026-11-11; simplified and Lerario removed
% at Kyle's request, relying on established thinkers). Sources: Axiom CFSG drew
% on 715 papers and books (research/01 s2, self-reported); Stacks Project tags
% (research/01 s3); Lean Atlas (Yanahama & Sannai, arXiv 2604.16347): review
% load cut by 94-99% for PNT+, Carleson and Brownian Motion (27% for XMSS, 69%
% for PhysLib; research/02 C). Thinkers, paraphrased: Thurston, "On proof and
% progress in mathematics", Bull. AMS 30(2), 1994; Lakatos, Proofs and
% Refutations, Cambridge UP, 1976 (proof as decomposing a conjecture into
% subconjectures or lemmas); Parnas, CACM 15(12), 1972; Wirth, CACM 14(4), 1971.
% The conclusion is scoped to the statement and proof criteria.
\begin{frame}{Lemmas as the unit of understanding}
  \centering
  \begin{tikzpicture}[x=1cm, y=1cm,
      colhd/.style={anchor=north west, inner sep=0pt, font=\sffamily\small\bfseries, text=black!75},
      body/.style={anchor=north west, text width=6.6cm, align=flush left, inner sep=0pt, font=\sffamily\footnotesize}]
    \useasboundingbox (0,0.1) rectangle (14.3,-5.9);
    \node[anchor=north west, text width=14.2cm, align=flush left, inner sep=0pt, font=\sffamily\small] at (0,0.05)
      {A proof is understood when a reader can walk its graph of lemmas: read each statement, and see why it
       follows from those beneath it, without opening a proof.};
    \draw[black!12, line width=0.5pt] (7.15,-1.05) -- (7.15,-4.55);
    \node[colhd] at (0,-1.1) {We already work this way};
    \node[body] at (0,-1.5) {%
      \begin{itemize}\setlength{\itemsep}{2pt}
        \item We cite results we have never read: the classification of finite simple groups rests on 715 papers and books.
        \item Reference works are graphs of lemmas: the Stacks Project is cited by tag.
        \item Many results in applied math and theoretical machine learning are not novel at the level of the lemma, but at the level of their \emph{assembly}.
      \end{itemize}};
    \node[colhd] at (7.5,-1.1) {Follows established thought};
    \node[body] at (7.5,-1.6) {%
      \begin{itemize}\setlength{\itemsep}{5pt}
        \item \textbf{Thurston} (1994): progress in mathematics is measured by human understanding, not by formal proof alone.
        \item \textbf{Lakatos} (1976): a proof decomposes a conjecture into subconjectures, or lemmas.
        \item \textbf{Parnas} (1972), \textbf{Wirth} (1971): use a component through its interface; refine one step at a time.
      \end{itemize}};
    \node[keybox, anchor=south, text width=12.6cm] at (7.15,-5.85)
      {Principled blueprints, statements aligned lemma by lemma, and generated formal proofs:\\
       enough for both understanding and integrity.};
  \end{tikzpicture}
  \slidesources{Axiom, CFSG (Oct 2026) and Stacks (Nov 2026) \textbullet{} W.~P.~Thurston, \emph{On proof and progress in mathematics}, Bull.\ AMS 30(2), 1994 \textbullet{} I.~Lakatos, \emph{Proofs and Refutations}, Cambridge UP, 1976 \textbullet{} D.~L.~Parnas, CACM 15(12), 1972 \textbullet{} N.~Wirth, CACM 14(4), 1971}
\end{frame}
